% ---
% id: protocol-compiler/rem:engine-admission
% kind: rem
% title: "Relation to the consuming engine's admission conditions"
% summary: >-
%   The consuming engine decides the same five exact conditions as this
%   module, and decides them alike. The declarations admitted here are a
%   subset of those the engine admits: D(i) and D(iii) are stronger here and
%   check V7 is one-sided.
% prominence: standard
% status: accepted
% uses: [def:exact, def:target]
% mentions: [foundations/req:consumer-decides-exactly, meat-core/def:admissible-declaration, def:validator]
% assumes: []
% refs: []
% provenance: newly-derived
% lean:
% ---

\begin{remark}\label{rem:engine-admission}
The five conditions of Definition~\ref{def:exact} are also the consuming
engine's admission conditions, which that engine decides exactly, allowing
no nonzero \tref[defects]{defect}
(\xref{foundations/req:consumer-decides-exactly}), so the two modules
decide \emph{those five conditions} alike. \emph{The agreement is on the five
conditions and not on the declarations they are conditions of}: this module
additionally imposes D(i)'s clause on $\idxm|_{\SecSet}$
(Definition~\ref{def:target}), for which that engine's definition of an
admissible declaration (\xref{meat-core/def:admissible-declaration}) has no
counterpart, and D(iii)'s Hermiticity clause, which that definition leaves to
the typing of the tuple;
and its feasibility check V7 (Definition~\ref{def:validator}) is one-sided:
it may refuse a declaration whose feasible set is non-empty, for want of a
verified witness. So the declarations admitted
here are a \emph{subset} of those that engine admits, and where the two
differ this module refuses declarations that engine would admit, never the
reverse.
\end{remark}
